% Generated from locale/tr/content/turing-machines/undecidability/representing-tms.tex; do not hand-edit.


\olfileid[tr]{tur}{und}{rep}
\olsection{Turing Makinelerinin Gösterilmesi}

\begin{explain}
Turing makinelerini ve davranışlarını birinci dereceden mantığın
!!a{sentence} ile göstermek için uygun bir dil tanımlamamız gerekir. Dil iki
bölümden oluşur: Makinenin konfigürasyonlarını betimleyen !!{predicate}s ve
çalıştırma adımlarını (``anları'') ve şeritteki konumları numaralandıran
ifadeler.

Her ikisi de iki yerli iki tür !!{predicate} tanıtırız: Her $q$ durumu
için !!a{predicate}~$\Obj Q_q$ ve her $\sigma$ şerit simgesi için
!!a{predicate}~$\Obj S_\sigma$. Birinciler $M$'nin durumunu ve şerit kafasının
konumunu, ikinciler şeridin içeriğini betimlememizi sağlar.

Şerit kafasının konumlarını ve yerine getirilen adımların sayısını ifade etmek
için sayıları ifade etmenin bir yoluna gereksinim duyarız. Bunun için
!!a{constant}~$\Obj 0$ ve ardıl fonksiyonu olan $1$ yerli $\prime$
fonksiyonunu kullanırız. Uzlaşım gereği bu fonksiyon argümanından \emph{sonra}
yazılır (ve parantezleri yazmayız).

Her $n$ sayısı için onu $\Lang L_M$ içinde gösteren bir kanonik terim
$\num{n}$, yani $n$'nin \emph{sayısalı} vardır. $\num{0}$, $\Obj 0$;
$\num{1}$, $\Obj 0'$; $\num{2}$, $\Obj 0''$ vb. biçimindedir. Daha biçimsel
olarak:
\begin{align*}
\num{0} & = \Obj 0 \\
\num{n+1} &= \num{n}'
\end{align*}
$\num{0}$ terimi, yani $\Obj 0$, şeritteki en soldaki konumun yanı sıra ilk
çalıştırma adımından önceki zamanı (başlangıç konfigürasyonunu) adlandırır.
$\num{1}$ terimi, yani $\Obj 0'$, en soldaki karenin sağındaki kareyi ve ilk
çalıştırma adımından sonraki zamanı adlandırır; bu böyle sürer.

Hem şerit konumlarının sıralamasını (``solunda'' anlamına geldiğinde) hem de
çalıştırma adımlarının sıralamasını (``önce'' anlamına geldiğinde) ifade etmek
için !!a{predicate}~$<$ da tanıtırız.

Dili kurduktan sonra $!T(M,w)$'nin ``aksiyomlarını'', yani birlikte ele
alındıklarında $M$'nin $w$ girdisinde çalıştırıldığındaki davranışını
betimleyen !!{sentence}s listeleriz. $\Obj 0$, $\prime$ ve $<$ üzerine
koşullar koyan !!{sentence}s, girdi konfigürasyonunu betimleyen !!{sentence}s
ve $M$ belirli bir yönergeyi yerine getirdikten sonraki konfigürasyonunu
betimleyen !!{sentence}s bulunacaktır.
\end{explain}

\begin{defn}
  \ollabel{defn:tm-descr}
Bir $M=\tuple{Q,\Sigma,q_0,\delta}$ Turing makinesi verildiğinde
$\Lang L_M$ dili şunlardan oluşur:
\begin{enumerate}
\item Her $q\in Q$ durumu için iki yerli bir !!{predicate}
  $\Obj Q_q(x,y)$. Sezgisel olarak $\Obj Q_q(\num{m},\num{n})$, ``$n$ adım
  sonra $M$, $m$. kareyi tararken $q$ durumundadır'' ifadesini dile getirir.
\item Her $\sigma\in\Sigma$ simgesi için iki yerli bir !!{predicate}
  $\Obj S_\sigma(x,y)$. Sezgisel olarak
  $\Obj S_\sigma(\num{m},\num{n})$, ``$n$ adım sonra $m$. kare $\sigma$
  simgesini içerir'' ifadesini dile getirir.
\item !!{constant} $\Obj 0$
\item Bir yerli !!{function} $\prime$
\item İki yerli !!{predicate} $<$
\end{enumerate}
\end{defn}

Turing makinesi~$M$'nin $w=\sigma_{i_1}\dots\sigma_{i_k}$ girdisindeki
işleyişini betimleyen !!{sentence}s şunlardır:
\begin{enumerate}
\item Sayıları ve~$<$ bağıntısını betimleyen aksiyomlar:
\begin{enumerate}
\item Her sayının ardılından küçük olduğunu söyleyen !!^a{sentence}:
\[
\lforall[x][x < x']
\]
\item $<$ bağıntısının geçişli olmasını sağlayan !!^a{sentence}:
\[
\lforall[x][\lforall[y][\lforall[z][
      ((x < y \land y < z) \lif x < z)]]]
\]
\end{enumerate}
\item Girdi konfigürasyonunu betimleyen aksiyomlar:
\begin{enumerate}
\item $0$ adım sonra, yani makine başlamadan önce, $M$ ilk durum~$q_0$'da ve
  $1$. kareyi taramaktadır:
\[
\Obj Q_{q_0}(\num{1}, \num{0})
\]
\item İlk $k+1$ kare $\TMendtape$, $\sigma_{i_1}$, \dots,
  $\sigma_{i_k}$ simgelerini içerir:
\[
\Obj S_\TMendtape(\num{0}, \num{0}) \land
\Obj S_{\sigma_{i_1}}(\num{1}, \num{0}) \land
\dots \land
\Obj S_{\sigma_{i_k}}(\num{k}, \num{0})
\]
\item Bunun dışında şerit boştur:
\[
\lforall[x][(\num{k} < x \lif \Obj S_\TMblank(x, \num{0}))]
\]
\end{enumerate}
\item Bir konfigürasyondan sonrakine geçişi betimleyen aksiyomlar:

Aşağıdakiler için $!A(x,y)$, $\sigma\in\Sigma$ olmak üzere
\[
\lforall[z][
  (((z < x \lor x < z) \land \Obj S_\sigma(z, y))
  \lif \Obj S_\sigma(z, y'))]
\]
biçimindeki bütün !!{sentence}s bağlacı olsun. $!A(\num{m},\num{n})$'yi,
``$m$. kare dışında şerit $n+1$ adım sonrasında $n$ adım sonrasıyla aynıdır''
ifadesini dile getirmek için kullanırız.
\begin{enumerate}
\item \ollabel{rep-right} Her
$\delta(q_i,\sigma)=\tuple{q_j,\sigma',\TMright}$ yönergesi için şu
!!{sentence}:
\begin{align*}
& \lforall[x][\lforall[y][(
   (\Obj Q_{q_i}(x, y) \land \Obj S_{\sigma}(x, y)) \lif {}]] \\
&\qquad   (\Obj Q_{q_j}(x', y') \land \Obj S_{\sigma'}(x, y') \land
!A(x, y)))
\end{align*}
Bu, makine $y$ adım sonra $\sigma$ simgesini içeren $x$. kareyi tararken
$q_i$ durumundaysa $y+1$ adım sonra $x+1$. kareyi taradığını, $q_j$ durumunda
olduğunu, $x$. karenin artık $\sigma'$ içerdiğini ve $x$ dışındaki her
karenin $y$ adım sonra içerdiği simgeyi içerdiğini söyler.

\item \ollabel{rep-left} Her
$\delta(q_i,\sigma)=\tuple{q_j,\sigma',\TMleft}$ yönergesi için şu
!!{sentence}:
\begin{align*}
& \lforall[x][\lforall[y][
    ((\Obj Q_{q_i}(x', y) \land \Obj S_{\sigma}(x', y)) \lif {}]]\\
& \qquad   (\Obj Q_{q_j}(x, y') \land \Obj S_{\sigma'}(x', y') \land
!A(x', y))) \land {}\\
& \lforall[y][((\Obj Q_{q_i}(\num{0}, y) \land \Obj S_{\sigma}(\num{0},
    y)) \lif {}]\\
& \qquad (\Obj Q_{q_j}(\num{0}, y') \land \Obj S_{\sigma'}(\num{0},
  y') \land !A(\num{0}, y)))
\end{align*}
Bunun nasıl çalıştığını düşünün: Artık ``$x$. kare taranıyorsa \dots'' diye
değil, ``$x+1$. kare taranıyorsa \dots'' diye başlarız. Sola hareket, sonraki
adımda makinenin $x$. kareyi taraması demektir. Fakat üzerine yazılan kare
$x+1$. karedir. Çıkarma ya da öncül fonksiyonumuz bulunmadığı için bunu böyle
yaparız.

$x+1$ biçimindeki sayıların $1$, $2$, \dots olduğuna dikkat edin; yani bu,
makinenin $0$. kareyi tarayıp sola gitmesi gereken durumu kapsamaz (elbette
bunu yapamaz, yalnızca yerinde kalır). Bu özel durum ikinci bağlaçla kapsanır:
Makine $y$ adım sonra $q_i$ durumunda $0$. kareyi tarıyor ve $0$. kare
$\sigma$ simgesini içeriyorsa $y+1$ adım sonra hâlâ $0$. kareyi taradığını,
artık $q_j$ durumunda olduğunu, $0$. karedeki simgenin $\sigma'$ olduğunu ve
$0$. kare dışındaki karelerin $y$ adım sonra içerdikleri simgeleri içermeyi
sürdürdüğünü söyler.
\item \ollabel{rep-stay} Her
$\delta(q_i,\sigma)=\tuple{q_j,\sigma',\TMstay}$ yönergesi için şu
!!{sentence}:
\begin{align*}
& \lforall[x][\lforall[y][(
   (\Obj Q_{q_i}(x, y) \land \Obj S_{\sigma}(x, y)) \lif {}]] \\
&\qquad   (\Obj Q_{q_j}(x, y') \land \Obj S_{\sigma'}(x, y') \land
!A(x, y)))
\end{align*}
\end{enumerate}
\end{enumerate}
Turing makinesi~$M$ ile $w$ girdisi için yukarıdaki bütün !!{sentence}s
bağlacı $!T(M,w)$ olsun.

$M$'nin sonunda durduğunu ifade etmek için ``belirli sayıda adımdan sonra
geçiş fonksiyonu tanımsız olacaktır'' diyen !!{sentence} bulmamız gerekir.
$X$, $\delta(q,\sigma)$ tanımsız olacak biçimdeki bütün
$\tuple{q,\sigma}$ çiftleri kümesi olsun. O zaman $!E(M,w)$ şu !!{sentence}
olsun:
\[
\lexists[x][\lexists[y][(\bigvee_{\tuple{q, \sigma} \in
      X}(\Obj Q_q(x, y) \land \Obj S_\sigma(x, y)))]]
\]

Özel olarak belirlenmiş bir durma durumu~$h$ olan bir Turing makinesi
kullanırsak daha da kolaydır: O zaman
\[
\lexists[x][\lexists[y][\Obj Q_h(x, y)]]
\]
!!{sentence}~$!E(M,w)$, makinenin sonunda durduğunu ifade eder.

\begin{prop}
\ollabel{prop:mlessk}
$m < k$ ise $!T(M, w) \Entails \num{m} < \num{k}$.
\end{prop}

\begin{proof}
Alıştırma.
\end{proof}

\begin{prob}
\olref[tur][und][rep]{prop:mlessk} önermesini kanıtlayın.
(İpucu: $k-m$ üzerine tümevarım kullanın.)
\end{prob}

